\documentclass[11pt]{article}
\usepackage[a4paper,margin=0.9in]{geometry}
\usepackage[T1]{fontenc}
\usepackage{helvet}
\renewcommand{\familydefault}{\sfdefault}
\usepackage{xcolor}
\usepackage{booktabs}
\usepackage{array}
\usepackage{colortbl}
\usepackage{float}
\usepackage{amsmath}
\usepackage{titlesec}
\usepackage{fancyhdr}
\usepackage{enumitem}
\usepackage{lastpage}
\usepackage[hidelinks]{hyperref}

\definecolor{bcnavy}{RGB}{11,31,58}
\definecolor{bcblue}{RGB}{32,74,135}
\definecolor{bcgold}{RGB}{176,141,87}
\definecolor{bcgrey}{RGB}{90,98,112}
\definecolor{bcrule}{RGB}{210,214,220}
\definecolor{bckeep}{RGB}{232,237,244}
\definecolor{bcact}{RGB}{247,241,228}
\definecolor{bcwarn}{RGB}{250,238,238}

\titleformat{\section}{\color{bcnavy}\large\bfseries}{\thesection.}{0.6em}{}
\titleformat{\subsection}{\color{bcblue}\normalsize\bfseries}{\thesubsection}{0.6em}{}
\pagestyle{fancy}\fancyhf{}
\lhead{\color{bcgrey}\small Bayesian Capital --- Book migration: five names to eight or nine}
\rfoot{\color{bcgrey}\small \thepage\ of \pageref{LastPage}}
\renewcommand{\headrulewidth}{0.4pt}

\newcommand{\pc}[1]{#1\%}

\begin{document}

\begin{center}
{\color{bcnavy}\LARGE\bfseries Migrating the book from five names to eight or nine}\\[4pt]
{\color{bcgrey}\large A justification, and the case against}\\[10pt]
{\color{bcgrey}\small 6 August 2026 \quad$\cdot$\quad Every figure verified against 5-minute bars \quad$\cdot$\quad
\texttt{premarket\_study/book\_verified.py}, \texttt{migration\_evidence.py}}
\end{center}

\vspace{4pt}
\noindent\fcolorbox{bcrule}{bckeep}{\parbox{\dimexpr\textwidth-2\fboxsep-2\fboxrule}{%
\textbf{Summary.} Migration costs roughly \pc{14} of planned return and buys a one-third
reduction in the book's exposure to a single theme. The deployed five plan at \pc{73.6}; the
eight-name book at \pc{59.7} and the nine-name at \pc{56.9}--\pc{59.1}. In exchange, average
pairwise correlation falls from 0.489 to 0.28--0.34, beta to the AI factor from 0.60 to
0.45--0.49, and the loss in the sample's worst theme drawdown from \pc{-39.0} to \pc{-29.1}.
\textbf{This is not a return improvement and should not be sold as one.} It is the purchase of
insurance at a stated price. Whether it is worth paying depends on a judgement about
concentration that the data informs but does not settle.}}

\section{What is being proposed}

The deployed book is five names --- RKLB, TSM, VST, VRT and MU --- each taking a fifth of the
capital and running two sleeves, Bayes and OU, at 50/50. The proposal is to widen it:

\begin{itemize}[leftmargin=1.4em,itemsep=2pt]
  \item \textbf{add NVDA and AVGO}, on a patient rule --- premium 4--6\%, 200-day stop,
        residual OU sigma with the buffer re-expressed in residual units;
  \item \textbf{add GM, VLO and CF}, the diversifier work's recommendation, on their fitted
        vectors unchanged;
  \item \textbf{optionally drop VST, or VLO, or both}, giving books of eight, nine or ten names.
\end{itemize}

Every name is equally weighted. Capital is allocated once and each name compounds
independently thereafter --- the \emph{held} construction. Rebalancing between names is not
merely avoided but forbidden: it feeds capital into falling names and sells winners, and when
it was done by accident it cost 16.3pp and inflated measured drawdown from \pc{14.9} to
\pc{29.7}.

\section{The basis, and why this document can be trusted more than its predecessors}

Every figure here is computed on \textbf{verified fills}. The model books a same-day round trip
only where 5-minute bars prove the low preceded the high. This matters more than any other
methodological choice in the study: the workbook's own cells report \pc{221} for VST and
\pc{438} for RKLB, and essentially the entire drop to the figures below is same-day fill
verification, not modelling disagreement.

Three specific improvements over earlier versions of this analysis:

\begin{enumerate}[leftmargin=1.4em,itemsep=3pt]
  \item \textbf{Nothing is quoted.} Earlier drafts mixed computed figures with numbers carried
        from previous sessions. Every row below is recomputed here from validated workbooks and
        5-minute files on one basis.
  \item \textbf{The at-open floor is no longer used for ranking.} An earlier draft ranked names
        at the at-open floor, which is a hard lower bound but a poor ordering device: measured
        against verified fills its error runs from $+10.1$pp (MU) to $+88.3$pp (RKLB), and it
        reverses the order of TSM and VRT. Conclusions drawn from it --- including the original
        case for dropping VST --- did not survive.
  \item \textbf{Planning figures use each name's own measured haircut.} The deployed names carry
        the 2.1pp out-of-sample allowance measured for them. GM, VLO and CF carry the
        unseen-window figures from the diversifier study, whose measured gap is far larger ---
        VLO gives back 29pp, GM 17pp. Applying 2.1pp to those three, as an earlier draft did,
        overstated each by roughly 20pp.
\end{enumerate}

\section{The names}

\begin{table}[H]\centering\small
\begin{tabular}{llrrrrrrr}
\toprule
\textbf{Name} & \textbf{Rule} & \textbf{Potential} & \textbf{Fitted} & \textbf{Tested} &
\textbf{Worst half} & \textbf{Planned} & \textbf{Buys/yr} & \textbf{MaxDD} \\
\midrule
RKLB & deployed    & 148.1\% & 162.5\% & 135.1\% & \textbf{135.1\%} & 146.0\% & 118 & 46.9\% \\
MU   & deployed    &  61.5\% &   7.1\% & 139.0\% & 7.1\%            &  59.4\% &  57 & 34.1\% \\
VST  & deployed    &  56.7\% &  53.4\% &  59.9\% & \textbf{53.4\%}  &  54.6\% &  79 & 38.9\% \\
VRT  & deployed    &  56.4\% &  -5.2\% & 152.2\% & -5.2\%           &  54.3\% &  74 & 62.7\% \\
TSM  & deployed    &  56.0\% &  34.3\% &  79.9\% & 34.3\%           &  53.9\% &  92 & 27.8\% \\
\midrule
AVGO & 4--6\%/200d &  48.7\% &  19.9\% &  82.5\% & 19.9\%           &  46.6\% &  19 & 38.7\% \\
NVDA & 4--6\%/200d &  45.2\% &  45.9\% &  44.5\% & \textbf{44.5\%}  &  43.1\% &  15 & 27.1\% \\
\midrule
VLO  & deployed    &  64.2\% &  11.4\% & 137.9\% & 11.4\%           &  35.0\% &  35 & 19.8\% \\
GM   & deployed    &  50.3\% &  44.0\% &  56.5\% & \textbf{44.0\%}  &  33.0\% &  61 & \textbf{14.6\%} \\
CF   & deployed    &  46.9\% &  39.2\% &  54.7\% & 39.2\%           &  41.0\% &  27 & 20.5\% \\
\bottomrule
\end{tabular}
\end{table}

\noindent Two observations that the ordering by \emph{potential} conceals.

\textbf{The incumbents are less reliable than they look.} Ranked by worst half --- the lower of
the two window returns, which is what a name is worth if its good window turns out to be the
fitted one --- the order is RKLB \pc{135}, VST \pc{53}, NVDA \pc{44}, GM \pc{44}, CF \pc{39},
TSM \pc{34}, AVGO \pc{20}, VLO \pc{11}, MU \pc{7}, VRT \pc{-5}. Three of the five deployed names
sit in the bottom four. VRT returns \pc{56.4} on halves of $-5.2$ and $152.2$; GM returns
\pc{50.3} on halves of $44.0$ and $56.5$. Those are not the same kind of \pc{50}.

\textbf{The published figures were optimistic.} Recomputed here: VST \pc{56.7} against a
published \pc{62}, VRT \pc{56.4} against \pc{67}, RKLB \pc{148.1} against \pc{158}. TSM and MU
come in close. The five-name book's published \pc{79} plan should be read as \pc{73.6}.

\section{The books}

\begin{table}[H]\centering\small
\begin{tabular}{lrrrrrrr}
\toprule
\textbf{Book} & \textbf{Potential} & \textbf{Fitted} & \textbf{Tested} & \textbf{Planned} &
\textbf{Buys/yr} & \textbf{MaxDD} & \textbf{Corr} \\
\midrule
\rowcolor{bckeep}
five (deployed)      & 80.8\% & 51.9\% & 113.4\% & \textbf{73.6\%} & 419 & 39.3\% & 0.489 \\
eight ($-$VST $-$VLO)& 68.5\% & 44.5\% &  95.0\% & \textbf{59.7\%} & 462 & 30.3\% & 0.322 \\
nine ($-$VLO)        & 67.2\% & 45.5\% &  90.9\% & 59.1\%          & 540 & 30.8\% & 0.339 \\
nine ($-$VST)        & 68.0\% & 40.9\% &  98.7\% & 56.9\%          & 497 & \textbf{29.1\%} & \textbf{0.283} \\
ten (everything)     & 66.9\% & 42.1\% &  94.6\% & 56.7\%          & 576 & 29.7\% & 0.299 \\
\bottomrule
\end{tabular}
\end{table}

\noindent The cost of migration is \textbf{13.9pp} of planned return to eight names, \pc{16.7}
to nine. Trade activity rises from 419 to 462--576 buys a year.

\subsection*{Why the return falls, and why that is not an argument about merit}

RKLB plans at \pc{146}. Every other name in the universe plans below \pc{60}. Under equal
weighting the book return is close to the average of its names, so \emph{any} addition below the
book average dilutes it, however good the added name is. The five-name book's \pc{73.6} is not a
property of the strategy; it is a property of owning RKLB at a fifth of the capital.

This cuts both ways and should be said plainly: the same arithmetic means the five-name plan
rests on one name. RKLB contributes 29.2pp of the \pc{73.6}. If RKLB merely matched the other
four the five-name plan falls to \pc{55.5}; if RKLB earned nothing it falls to \pc{44.4}. The
eight-name book's \pc{59.7} is lower but rests on eight legs rather than one.

\section{What the migration actually buys}

\subsection*{Concentration}

\begin{table}[H]\centering\small
\begin{tabular}{lrrr}
\toprule
\textbf{Book} & \textbf{Avg pairwise corr} & \textbf{Beta to AI factor} & \textbf{$R^2$} \\
\midrule
five (deployed)       & 0.489 & 0.603 & 0.391 \\
eight ($-$VST $-$VLO) & 0.322 & 0.484 & 0.423 \\
nine ($-$VLO)         & 0.339 & 0.489 & 0.471 \\
nine ($-$VST)         & \textbf{0.283} & \textbf{0.452} & 0.427 \\
ten (everything)      & 0.299 & 0.461 & 0.473 \\
\bottomrule
\end{tabular}
\end{table}

The deployed book is one trade. TSM is semis, VRT data-centre cooling, VST data-centre power, MU
HBM memory; RKLB is not an AI name but carries beta 0.77 to the factor and 0.69 to TSM --- it
shares the risk appetite rather than the theme. Five-for-five exposed.

Note that $R^2$ barely moves while beta falls by a quarter. Widening the book scales total
exposure to the theme \emph{down}; it does not change what the exposure \emph{is}. Beta and the
stress windows below are the honest measures. $R^2$ is not, and a document that led with it
would be overstating the case.

\subsection*{Behaviour in the theme's worst drawdowns}

Windows located automatically as the AI factor's deepest peak-to-trough declines in the sample,
not chosen.

\begin{table}[H]\centering\small
\begin{tabular}{lrrr}
\toprule
& \textbf{Jan--Apr 2025} & \textbf{Jun--Jul 2026} & \textbf{May--Aug 2024} \\
& \textbf{factor $-$45.9\%} & \textbf{factor $-$27.9\%} & \textbf{factor $-$27.3\%} \\
& \emph{fitted half} & \emph{\textbf{tested half}} & \emph{fitted half} \\
\midrule
five (deployed)       & $-$39.0\% & $-$27.8\% & $-$5.2\% \\
eight ($-$VST $-$VLO) & $-$30.3\% & $-$20.3\% & $-$6.2\% \\
nine ($-$VLO)         & $-$30.8\% & $-$19.2\% & $-$5.4\% \\
nine ($-$VST)         & $-$29.1\% & $-$17.2\% & $-$5.6\% \\
ten (everything)      & $-$29.7\% & $-$16.5\% & $-$4.9\% \\
\bottomrule
\end{tabular}
\end{table}

\noindent\fcolorbox{bcrule}{bcact}{\parbox{\dimexpr\textwidth-2\fboxsep-2\fboxrule}{%
\textbf{A correction to the record.} The handover states that the AI factor rose in every
walk-forward test window, so the sample \emph{prices the premium for diversifying and never
tests the insurance}. That is true of the walk-forward folds. It is \textbf{not} true of the
half-sample split used throughout this document: Jun--Jul 2026 is a \pc{-27.9} factor drawdown
sitting squarely in the tested half, and the wider books lost 8--11pp less in it than the
deployed five. \emph{Caveat, stated rather than buried:} the diversifier vectors are full-sample
fits, so that window lies inside their fitted region too. It is not clean out-of-sample
evidence. It is like-for-like --- the deployed five are also full-sample fits on the same window
--- and it is the closest thing to a live test of the insurance the sample contains.}}

\section{The case against migration}

A justification that only argues one side is not a justification. The strongest objections:

\begin{enumerate}[leftmargin=1.4em,itemsep=4pt]
  \item \textbf{Roughly 14pp is a large, certain price for an uncertain benefit.} The return
        cost is arithmetic and will be paid. The diversification benefit is a claim about states
        of the world that have occurred three times in 2.3 years.
  \item \textbf{Two of the three stress windows are in the fitted half.} Only Jun--Jul 2026 is
        not, and even that sits inside the diversifiers' fitted region.
  \item \textbf{The added names' out-of-sample record is weak where it has been measured.}
        GM gives back 17pp from full fit to unseen windows and VLO 29pp. Their planning figures
        are \pc{33} and \pc{35} against full fits of \pc{50.3} and \pc{64.2}. The honest
        expectation for these names is a third to a half of what the backtest shows.
  \item \textbf{NVDA and AVGO carry an unmeasured haircut.} Their planning figures apply the
        deployed names' 2.1pp allowance, but nothing establishes a haircut at the 4--6\%/200-day
        setting. These are the two softest numbers in the document.
  \item \textbf{The strategy loses to buy-and-hold on the incumbents.} Over this sample the
        rule beats outright buy-and-hold on only 5 of 10 names. If the conviction is simply that
        these names rise, widening the book is not the first argument to have --- holding more
        stock is.
\end{enumerate}

\noindent However, that last objection inverts on inspection, and the inversion is the most
interesting single result in this document:

\begin{table}[H]\centering\small
\begin{tabular}{lrrl}
\toprule
\textbf{Name} & \textbf{Strategy} & \textbf{Buy-and-hold} & \textbf{Rule adds value?} \\
\midrule
RKLB & 148.1\% & 238.1\% & no \\
MU   &  61.5\% & 125.2\% & no \\
VRT  &  56.4\% &  66.0\% & no \\
AVGO &  48.7\% &  57.8\% & no \\
TSM  &  56.0\% &  57.0\% & no \\
\midrule
\rowcolor{bckeep} VLO  &  64.2\% & 28.0\% & \textbf{yes, $+$36pp} \\
\rowcolor{bckeep} CF   &  46.9\% & 16.2\% & \textbf{yes, $+$31pp} \\
\rowcolor{bckeep} GM   &  50.3\% & 32.5\% & \textbf{yes, $+$18pp} \\
\rowcolor{bckeep} VST  &  56.7\% & 39.6\% & \textbf{yes, $+$17pp} \\
\rowcolor{bckeep} NVDA &  45.2\% & 42.4\% & \textbf{yes, $+$3pp} \\
\bottomrule
\end{tabular}
\end{table}

\noindent The limit-order book \emph{loses} to simply holding the stock on four of five AI names
and \emph{beats} it on all three diversifiers. That is what a dip-buying, premium-selling rule
should do: it gives up upside in a runaway trend and earns in a choppy or range-bound one. The
regime work says the same thing from the other direction --- the deployed rule's
exposure-matched edge is $+$\pc{14.1} when the theme is weak and $-$\pc{17.3} when it is
running.

The implication is uncomfortable for the five-name book and favourable to migration: over this
sample the deployed five have been a leveraged bet on the AI theme executed through a rule that
is \emph{worse} than holding the theme outright. The names on which the strategy genuinely earns
its keep are the ones being proposed for addition.

\section{Recommendation}

\noindent\fcolorbox{bcrule}{bcact}{\parbox{\dimexpr\textwidth-2\fboxsep-2\fboxrule}{%
\textbf{Migrate to eight names} --- RKLB, TSM, VRT, MU, NVDA, AVGO, GM, CF --- at \pc{12.5}
each, for a planned \pc{59.7} on 462 buys a year, correlation 0.322 and beta 0.484. Retain VST
if a ninth slot is acceptable; the case for dropping it is the weakest decision in this
document.}}

\subsection*{On the two disputed names}

\textbf{VLO: drop.} Halves of \pc{11.4} and \pc{137.9} for a \pc{64.2} average is the
wide-fold-spread overfitting tell, and its measured out-of-sample give-back of 29pp is the
largest of any name. Dropping it raises the planned book \pc{56.7} $\to$ \pc{59.7} and costs
0.039 of correlation.

\textbf{VST: marginal, and previously mis-argued.} VST was originally dropped on the observation
that it was weakest at the at-open floor. That basis is disqualified. On verified fills VST is
indeed the lowest of the deployed five at \pc{56.7} --- but only just, against TSM \pc{56.0} and
VRT \pc{56.4} --- while being \emph{second of ten} by worst half at \pc{53.4}. It is the
steadiest earner in the deployed book, not the weakest one. Dropping it buys 1.1pp of planned
return and 0.017 of correlation, and removes the book's most reliable non-RKLB name. Either
choice is defensible; the eight-name recommendation above reflects a preference for fewer
positions rather than a finding against VST.

\subsection*{What would change this recommendation}

\begin{itemize}[leftmargin=1.4em,itemsep=2pt]
  \item A reversed walk-forward --- fit on the late sample, score the Jan--Apr 2025 crash ---
        would put the insurance genuinely out of sample. It remains the single most valuable
        outstanding test and has been offered before.
  \item A measured haircut for NVDA and AVGO at the 4--6\%/200-day setting.
  \item Any evidence that RKLB's \pc{146} is not repeatable would strengthen the case for
        migration sharply, since the five-name plan depends on it.
\end{itemize}

\vspace{6pt}
\noindent\fcolorbox{bcrule}{bcwarn}{\parbox{\dimexpr\textwidth-2\fboxsep-2\fboxrule}{%
\textbf{What this document does not establish.} That the eight-name book will return \pc{59.7}.
Planning figures are haircut backtests on 2.3 years of one market regime, and the interval
around every number here is far wider than the differences being discussed. The defensible claim
is comparative --- that the wider book carries roughly one-third less exposure to a single theme,
for roughly 14pp of expected return --- not that either level will be realised.}}

\end{document}
